\chapter{Methadone}
 
\section*{Definition and Overview}
Methadone is a long-acting, synthetic opioid agonist primarily used in the management of opioid dependence (opioid agonist therapy, or OAT) and for the treatment of severe, chronic pain refractory to other analgesics. Pharmacologically, methadone binds to mu-opioid receptors in the central nervous system, mimicking endogenous opioids. Due to its long half-life (typically 24 to 36 hours, and up to 50 hours in some individuals) and slow onset of action, methadone prevents withdrawal symptoms and reduces cravings for shorter-acting opioids (like heroin or oxycodone) without producing the rapid, intense euphoria associated with illicit use. In Australia, methadone is a Schedule 8 controlled drug, tightly regulated, and administered through specialized clinics and pharmacies to ensure safety and monitor adherence.
 
\section*{Epidemiology}
Methadone maintenance therapy is a cornerstone of harm reduction programs in Australia and worldwide. Approximately 50,000 Australians receive pharmacotherapy for opioid dependence on any given day, with methadone accounting for roughly 50--60\% of these cases (the remainder using buprenorphine). The population receiving methadone is predominantly male, with a median age in the early 40s. While highly effective, methadone is a common contributor to pharmaceutical opioid overdose presentations, often due to accidental ingestion by non-tolerant individuals or combination with other respiratory depressants.
 
\section*{Aetiology and Pathophysiology}
Methadone's therapeutic effects and clinical risks are explained by its complex pharmacokinetic and pharmacodynamic properties:
 
\begin{itemize}
    \item \textbf{Mu-opioid receptor agonism:} methadone is a full agonist at mu-opioid receptors, which inhibits adenylyl cyclase, closes voltage-gated calcium channels, and opens potassium channels, reducing neuronal excitability and pain transmission.
    \item \textbf{NMDA receptor antagonism:} unlike most opioids, methadone acts as an antagonist at N-methyl-D-aspartate (NMDA) receptors, which helps prevent opioid tolerance and is beneficial for neuropathic pain.
    \item \textbf{Long and variable half-life:} its prolonged clearance is due to extensive tissue binding. Slow dissociation from receptors prevents rapid fluctuations in plasma levels, stabilizing the patient's physiological state.
    \item \textbf{Cardiac conduction effects:} methadone blocks the hERG potassium channel, which can delay cardiac repolarisation and cause QT interval prolongation.
\end{itemize}
 
\section*{Clinical Features}
The clinical effects range from therapeutic stabilization to signs of acute toxicity and withdrawal:
 
\begin{itemize}
    \item Therapeutic effects: prevention of opioid withdrawal symptoms, suppression of drug cravings, and effective relief of chronic pain
    \item Common side effects: constipation, sweating (diaphoresis), dry mouth, drowsiness, weight gain, and sexual dysfunction
    \item Signs of acute toxicity or overdose: miosis (pinpoint pupils), central nervous system depression (stupor, coma), bradycardia, and severe respiratory depression (slow, shallow breathing)
    \item Cardiac features: palpitations or syncope related to QT prolongation and torsades de pointes
    \item Withdrawal features: anxiety, restlessness, sweating, lacrimation (tears), rhinorrhea (runny nose), dilated pupils, muscle aches, abdominal cramps, nausea, vomiting, and diarrhoea (typically developing 24--48 hours after the last dose)
\end{itemize}
 
\section*{Differential Diagnosis}
\begin{itemize}
    \item \textbf{Other sedative or opioid overdoses (e.g., fentanyl, heroin, benzodiazepines):} presents similarly with respiratory depression, requiring history and specific drug screens to differentiate.
    \item \textbf{Severe hypothyroidism or myxoedema coma:} presents with bradycardia and hypothermia, but lacks pinpoint pupils and history of drug exposure.
    \item \textbf{Adrenal crisis:} presenting with hypotension and nausea, but not accompanied by pinpoint pupils or respiratory depression.
    \item \textbf{Long QT syndrome of other causes:} congenital long QT or induced by other medications (like erythromycin or haloperidol), differentiated by medication history and ECG analysis.
\end{itemize}
 
\section*{When to Seek Medical Care}
Immediate medical evaluation is required if someone is showing signs of methadone toxicity or if an accidental ingestion is suspected.
 
\textbf{Call 000 (or local emergency services) for:}
\begin{itemize}
    \item Unresponsiveness, severe lethargy, or inability to wake up
    \item Extremely slow, shallow, or irregular breathing, or gurgling sounds
    \item Pale, cold, or blue skin, lips, or fingernails (cyanosis)
    \item Palpitations, dizziness, or fainting episodes (suggesting a cardiac arrhythmia)
    \item Accidental ingestion of methadone by a child or non-tolerant adult
\end{itemize}
 
\section*{Diagnosis and Investigations}
\begin{itemize}
    \item \textbf{Clinical examination:} assessing airway, breathing, circulation, level of consciousness, and pupillary size.
    \item \textbf{Electrocardiogram (ECG):} critical to monitor the QTc interval, especially when starting treatment or increasing doses.
    \item \textbf{Urine or blood toxicology:} to detect methadone and its primary metabolite (EDDP), and to screen for other co-ingested substances (especially alcohol or benzodiazepines).
    \item \textbf{Blood tests:} renal and liver function tests, and electrolytes (imbalances can worsen QT prolongation).
\end{itemize}
 
\section*{Management}
Treatment of acute toxicity requires securing the airway and reversing respiratory depression:
 
\begin{itemize}
    \item \textbf{Airway management:} ensuring a patent airway and providing mechanical ventilation or high-flow oxygen.
    \item \textbf{Naloxone therapy:} the specific opioid antagonist is administered. Because methadone is long-acting, a continuous naloxone infusion is often required, as the duration of action of naloxone (30--90 minutes) is much shorter than that of methadone.
    \item \textbf{Long-term maintenance:} structured methadone maintenance programs under close medical supervision, including daily dosing, supervised consumption at pharmacies, and regular clinical reviews.
    \item \textbf{ECG monitoring:} routine ECG checks to monitor QTc interval, especially if doses exceed 100 mg daily.
    \item \textbf{Psychosocial support:} counselling, case management, and peer support to address addiction.
\end{itemize}
 
\section*{Complications}
\begin{itemize}
    \item Life-threatening respiratory depression and hypoxic brain damage
    \item Prolongation of the QT interval, leading to torsades de pointes and sudden cardiac arrest
    \item Severe, chronic constipation and bowel obstruction
    \item Neonatal abstinence syndrome (NAS) in infants born to mothers on methadone
    \item Development of physical dependence and tolerance
\end{itemize}
 
\section*{Prognosis and Outcomes}
With prompt supportive care and naloxone, the prognosis for recovery from acute methadone overdose is excellent. For opioid dependence, methadone maintenance has a highly favorable prognosis: it is associated with a dramatic reduction in illicit drug use, criminality, transmission of blood-borne viruses (HIV, Hepatitis C), and all-cause mortality. Long-term outcomes are optimized when pharmacotherapy is combined with vocational, psychological, and social rehabilitation.
 
\section*{Prevention and Self-Care}
\begin{itemize}
    \item Take methadone strictly as prescribed, and never take extra doses to compensate for missed days without consulting your doctor
    \item Avoid combining methadone with alcohol, sleeping pills, or other sedatives, which increases the risk of fatal overdose
    \item Store methadone in a child-resistant container, locked securely out of reach of children and others
    \item Attend all scheduled medical appointments and pharmacy dosing sessions regularly
    \item Learn to recognize signs of overdose and ensure family or housemates know where naloxone is kept and how to use it
\end{itemize}
 
\section*{Sources and Further Reading}
The following reputable organisations provide further information for healthcare students and patients seeking reliable, current advice about methadone.
 
\begin{itemize}
    \item Healthdirect Australia. \textit{Methadone: uses, safety, and support for opioid dependence.} https://www.healthdirect.gov.au/
    \item Therapeutic Goods Administration (TGA). \textit{National guidelines for medication-assisted treatment of opioid dependence.} https://www.tga.gov.au/
    \item Mayo Clinic. \textit{Methadone (oral route): precautions and side effects.} https://www.mayoclinic.org/
    \item MSD Manual. \textit{Opioid toxicity and withdrawal: methadone management.} https://www.msdmanuals.com/
    \item World Health Organization. \textit{Guidelines for the psychosocially assisted pharmacological treatment of opioid dependence.} https://www.who.int/
\end{itemize}
